% Slides for ApacheCon US 2015
%
% Phil Steitz <phil.steitz@gmail.com>
% April 14, 2015

\documentclass[14pt,mathserif]{beamer}
\usetheme{Boadilla}

\setbeamertemplate{footline}
{
  \leavevmode%
  \hbox{%
    \begin{beamercolorbox}[wd=.333333\paperwidth,ht=2.25ex,dp=1ex,center]{author in head/foot}%
    \usebeamerfont{author in head/foot}\insertshortinstitute
    \end{beamercolorbox}%
    \begin{beamercolorbox}[wd=.333333\paperwidth,ht=2.25ex,dp=1ex,center]{title in head/foot}%
      \usebeamerfont{title in head/foot}\insertshorttitle
    \end{beamercolorbox}%
    \begin{beamercolorbox}[wd=.333333\paperwidth,ht=2.25ex,dp=1ex,right]{date in head/foot}%
      \usebeamerfont{date in head/foot}\insertshortdate{}\hspace*{2em}
      \insertframenumber{} / \inserttotalframenumber\hspace*{2ex}
    \end{beamercolorbox}}%
  \vskip0pt%
}

%\useoutertheme[footline=authortitle]{miniframes}

\usepackage[T1]{fontenc}
\usepackage{lmodern}
\usepackage{rotating} % for defining \schwa
\usepackage{color}
\usepackage{minted}
\usepackage{longtable}

\newcommand{\fmgood}[1]{\color{green}{#1}}
\newcommand{\fmmedium}[1]{\color{orange}{#1}}
\newcommand{\fmbad}[1]{\color{red}{#1}}

\newcommand{\newauthor}[2]{
  \parbox{0.26\textwidth}{
    \texorpdfstring
      {
        \centering
        #1 \\
        {\scriptsize{\urlstyle{same}\url{#2}\urlstyle{tt}}}
      }
      {#1}
  }
}

\definecolor{links}{HTML}{2A1B81}
\hypersetup{colorlinks,linkcolor=cyan,urlcolor=magenta}

\newcommand{\schwa}{\raisebox{1ex}{\begin{turn}{180}e\end{turn}}}

\AtBeginSection[]{
  \begin{frame}
  \vfill
  \centering
  \begin{beamercolorbox}[sep=8pt,center,shadow=true,rounded=true]{title}
    \usebeamerfont{title}\secname\par%
  \end{beamercolorbox}
  \vfill
  \end{frame}
}

%gets rid of bottom navigation bars
%\setbeamertemplate{footline}[page number]{}

%gets rid of navigation symbols
\setbeamertemplate{navigation symbols}{}

\title{Programming Math in Java}
\subtitle{Lessons from Apache Commons Math}
\author{
  \newauthor{Phil Steitz}{psteitz@apache.org}
}
\institute[Apachecon North America 2015]{Apachecon North America 2015}
\date{April 14, 2015}
\begin{document}

%----------- titlepage ----------------------------------------------%
{
\begin{frame} %[plain]
  \titlepage
\end{frame}
}

%----------- slide --------------------------------------------------%
\begin{frame}
  \frametitle{Agenda}
\begin{itemize}
  \item Commons Math Overview
  \item Examples and problems
  \item Getting involved
\end{itemize}

\end{frame}

%----------- slide --------------------------------------------------%
\begin{frame}
  \frametitle{Goals}

\begin{itemize}
  \item 100\% pure, self-contained Java
  \item Provide simple solutions to common problems
  \item Implement standard, well-documented algorithms
  \item Balance ease of use and good design
  \item Stay well-tested and maintained
  \item Maintain a strong community
\end{itemize}

\end{frame}

%----------- slide --------------------------------------------------%
\begin{frame}
  \frametitle{Message from our sponsors...}

\begin{columns}
\begin{column}{.5\textwidth}
Come join us! \\
 Code is interesting! \\
 Math is interesting! \\
 Bugs are interesting! \\
People are interesting! \\
Come join us! \\
\end{column}
\begin{column}{.5\textwidth}\raggedleft
\includegraphics[width=6cm]{commons-logo.png}\\~\\
\includegraphics[width=5cm]{math.jpg}
\end{column}
\end{columns}
\\
\url{http://commons.apache.org/math}
\end{frame}

%--------------------------------------------------------------------%
\section[Overview]{Quick Tour}
%----------- slide --------------------------------------------------%
\begin{frame}
  \frametitle{}

\begin{description}
  \item[Building Blocks] util, special, complex packages
\end{description}

\begin{small}
\begin{itemize}
  \item FastMath (pure java, portable replacement for java.lang.Math / StrictMath)
  \item Dfp (arbitrary precision arithmetic)
  \item Continued fractions
  \item Special functions (Beta, Gamma, Bessel, Erf)
  \item Array utilities (norms, normalization, filling, etc)
  \item Basic combinatorics (binomial coefficients, factorials, Stirling numbers..)
  \item Basic integer arithmetic (checked operations, gcd, lcm, primality...)
  \item Complex numbers
\end{itemize}
\end{small}
\end{frame}

%----------- slide --------------------------------------------------%
\begin{frame}
  \frametitle{}
  
\begin{description}
  \item[Linear Algebra] linear package 
\end{description}

\begin{small}
\begin{itemize}
  \item Vector and Matrix operations (dense, sparse, field)
  \item Decompositions (LU, QR, Cholesky, SVD, Eigenvalue)
  \item Solvers
\end{itemize}
\end{small}

\begin{description}
  \item[Basic Numerical Analysis] analysis package
\end{description}

\begin{small}
\begin{itemize}
  \item Automatic Differentiation
  \item Numerical Integration
  \item Interpolation
  \item Root finders
\end{itemize}
\end{small}


\end{frame}

%----------- slide --------------------------------------------------%
\begin{frame}
  \frametitle{}

\begin{description}
  \item[Probability and Statistics] distributions, stat and random packages
\end{description}

\begin{small}
\begin{itemize}
  \item Probability distributions
  \item Random number generators (Well, ISAAC, Mersenne twister...)
  \item Random data generators (samplers, random vectors, empirical density)
  \item Univariate statistics (storeless and in-memory)
  \item Regression (storeless bivariate, matrix-based and storeless multivariate)
  \item Correlation
  \item Inference (T-, G-, ChiSquare, Kolmogorov-Smirnov tests...)  
\end{itemize}
\end{small}
\end{frame}

%----------- slide --------------------------------------------------%
\begin{frame}
  \frametitle{}
  
\begin{description}
  \item[Optimization and Geometry] optim, fitting, ode, transform, geometry 
\end{description}

 \begin{small}
 \begin{itemize}
  \item Fitting (non-linear least-squares, curve fitting)
  \item Ordinary differential equations (ODE)
  \item Linear/Non-linear optimization
  \item Computational Geometry
  \item Fast Fourier transform
 \end{itemize}
\end{small}
\end{frame}

%----------- slide --------------------------------------------------%
\begin{frame}
  \frametitle{}

\begin{description}
  \item[AI / Machine learning] genetics, filter and ml packages
\end{description}

\begin{small}
\begin{itemize}
  \item Genetic algorithms
  \item Neural networks
  \item Clustering
  \item Kalman filter
\end{itemize}
\end{small}
\end{frame}

%--------------------------------------------------------------------%
\section[Examples]{Examples and Challenges}
%----------- slide --------------------------------------------------%
\begin{frame}[fragile]
  \frametitle{Example 1 - Root finding}

\begin{small}
Find all roots of the polynomial \[p(x) = x^5 + 4x^3 + x^2 + 4 = (x+1)(x^2-x+1)(x^2+4)\]
\end{small}
\begin{minted}[xleftmargin=10pt,fontsize=\footnotesize,linenos=false]{java}

double[] coefficients = { 4.0, 0.0, 1.0, 4.0, 0.0, 1.0 };
LaguerreSolver solver = new LaguerreSolver();
Complex[] result = solver.solveAllComplex(coefficients, 0);
\end{minted}
\end{frame}

%----------- slide --------------------------------------------------%
\begin{frame}[fragile]
  \frametitle{Example - Root finding (cont)}

How about this one? \[p(x) = \prod\limits_{i=1}^{50}{15,000(i + ix + ix^2 + ix^3)}\]
As of 3.4.1, this will appear to hang.\\
Two problems:
\begin{enumerate}
  \item Complex iterates "escape to NaN"
  \item Default max function evaluations is set to Integer.MAX\_VALUE
\end{enumerate}
\end{frame}

%----------- slide --------------------------------------------------%
\begin{frame}[fragile]
  \frametitle{Example 1 - Root finding (cont)}

\begin{small}  
 First problem:  escape to NaN
\begin{itemize}
  \item Long and painful history debating C99x Annex G
  \item Desire: clear, self-contained, open documentation
  \item Desire: simple, performant code
  \item Result: use standard definitions and let NaNs propagate
\end{itemize}
Second problem: hanging
\begin{itemize}
  \item Ongoing debate about what to make defaults, whether to have defaults at all...
  \item Ease of use vs foot-shooting potential
  \item OO vs procedural / struct / primitives design
\end{itemize}
\end{small}
\end{frame} 
%----------- slide --------------------------------------------------%
\begin{frame}[fragile]
  \frametitle{Example 2 - Kolmogorov-Smirnov Test}

\begin{small}
Problem:  Given two samples \(S_1\) and \(S_2\), are they drawn from the same underlying
probability distribution?
\end{small}
\\
\\
\begin{minted}[xleftmargin=10pt,fontsize=\footnotesize,linenos=false]{java}

// load sample data into double[] arrays
double[] s1 = ... 
double[] s2 = ...
KolmogorovSmirnovTest test = new KolmogorovSmirnovTest();
// Compute p-value for test
double pValue = test.kolmogorovSmirnovTest(s1, s2);

\end{minted}
\end{frame}

%----------- slide --------------------------------------------------%
\begin{frame}
  \frametitle{K-S Test Implementation Challenges}
\begin{small}
Test statistic \( = D_{m,n} = \) maximum difference between the empirical
distribution functions of the two samples. \\
 \begin{itemize}
  \item The distribution of \(D_{m,n}\) is asymptotically an easy-to-compute distribution
  \item For very small m, n, the distribution can be computed exactly, but expensively
  \item What to do for mid-size samples (m, n between 100 and 10000)?
  \item Should this distribution be in the distributions package?
  \item Reference data / correctness very hard to verify
 \end{itemize}
\end{small}
\end{frame}
%----------- slide --------------------------------------------------%
\begin{frame}[fragile]
  \frametitle{Example 3 - Multivariate optimization}

\begin{small}
Find the minimum value of \(100(y-x^2)^2 + (1-x)^2\), starting at the point \((-1, 1)\)
\end{small}
\\
\begin{minted}[xleftmargin=10pt,fontsize=\scriptsize,linenos=false]{java}

MultivariateFunction f = new MultivariateFunction() {
    public double value(double[] x) {
        double xs = x[0] * x[0];
        return (100 * FastMath.pow((x[1] -  xs), 2)) +
           FastMath.pow(1 - x[0], 2);
    }
};
PowellOptimizer optim = new PowellOptimizer(1E-13, 1e-40);
PointValuePair result = optim.optimize(new MaxEval(1000),
                        new ObjectiveFunction(f),
                        GoalType.MINIMIZE,
                        new InitialGuess(new double[] {-1, 1}));
\end{minted}
\\
\begin{small}
Works OK in this case, but...
\end{small}
\end{frame}

%----------- slide --------------------------------------------------%
\begin{frame}[fragile]
  \frametitle{Example 3 - Multivariate optimization (cont)}
  
%-- ADD: Difficulty in stabilizing optimization API and how to steer users --%
\end{frame}
%----------- slide --------------------------------------------------%
\begin{frame}[fragile]
  \frametitle{Example 4 - Aggregated statistics}

\begin{small}
Problem: Aggregate summary statistics across multiple processes
\end{small}
\\
\begin{minted}[xleftmargin=10pt,fontsize=\scriptsize,linenos=false]{java}
// Create a collection of stats collectors
Collection<SummaryStatistics> aggregate =
     new ArrayList<SummaryStatistics>();
     
// Add a collector, hand it out to some process
// and gather some data...
SummaryStatistics procStats = new SummaryStatistics();
aggregate.add(procStats);
procStats.addValue(wildValue);
...
// Aggregate results
StatisticalSummary aggregatedStats =
     AggregateSummaryStatistics.aggregate(aggregate);

\end{minted}
\\
\begin{small}
\begin{itemize}
\item SummaryStatistics instances handle large data streams
\item StatisticalSummary is reporting interface
\end{itemize}
Question: how to distribute instances / workload?
\end{small}
\end{frame}

%----------- slide --------------------------------------------------%
\begin{frame}
  \frametitle{Threading and workload management}

\\
\begin{enumerate}
\item Should we implement multi-threaded algorithms?
\item How can we be more friendly for Hadoop / other distributed compute environments?
\end{enumerate}
\\
\\
So far (as of 3.4.1), we have not done 1. \\
So far, answer to 2 has been "leave it to users." \\
Could be both answers are wrong...
\end{frame}

%----------- slide --------------------------------------------------%
\begin{frame}[fragile]
  \frametitle{Example 5 - Genetic algorithm}

\begin{small}
Another example where concurrent execution would be natural...
\end{small}
\begin{minted}[xleftmargin=10pt,fontsize=\scriptsize,linenos=false]{java}
// initialize a new genetic algorithm
GeneticAlgorithm ga = new GeneticAlgorithm(
            new OnePointCrossover<Integer>(),
            CROSSOVER_RATE, 
            new BinaryMutation(),
            MUTATION_RATE,
            new TournamentSelection(TOURNAMENT_ARITY)
);

// initial population 
Population initial = getPopulation(); 

// stopping conditions
StoppingCondition stopCond = new FixedGenerationCount(NUM_GENERATIONS);

// run the algorithm
Population finalPopulation = ga.evolve(initial, stopCond);

// best chromosome from the final population
Chromosome bestFinal = finalPopulation.getFittestChromosome();

\end{minted}
 
\begin{small}
Question: how to distribute "evolve" workload?
\end{small}
\end{frame}

%----------- slide --------------------------------------------------%
\begin{frame}[fragile]
  \frametitle{Example 6 - OLS regression}

\begin{small}
Three ways to do it
\begin{description}
  \item[OLSMultipleRegression (stat.regression)] in-memory OLS
  \item[MillerUpdatingRegression (stat.regression)] streaming OLS
  \item[LevenbergMarquardtOptimizer (fitting.leastsquares)] in-memory
\end{description}
\begin{enumerate}
\item How to make sure users discover the right solution for them?
\item How much dependency / reuse / common structure to force?
\end{enumerate}
\end{small}
\end{frame}


%----------- slide --------------------------------------------------%
\begin{frame}[fragile]
  \frametitle{Example 6 - OLS regression (cont)}

\begin{small}
Simplest, most common:
\begin{minted}[xleftmargin=10pt,fontsize=\scriptsize,linenos=false]{java}
OLSMultipleLinearRegression regression = new OLSMultipleLinearRegression();
double[] y = new double[]{11.0, 12.0, 13.0, 14.0, 15.0, 16.0};
double[] x = new double[6][];
x[0] = new double[]{0, 0, 0, 0, 0};
x[1] = new double[]{2.0, 0, 0, 0, 0};
...
x[5] = new double[]{0, 0, 0, 0, 6.0};          
regression.newSample(y, x);
double[] beta = regression.estimateRegressionParameters();       
double[] residuals = regression.estimateResiduals();
double rSquared = regression.calculateRSquared();
double sigma = regression.estimateRegressionStandardError();
\end{minted}
\end{small}
Question:  What happens if design matrix is singular?
\end{frame}

%----------- slide --------------------------------------------------%
\begin{frame}[fragile]
  \frametitle{Example 6 - OLS regression (cont)}

\begin{small}
\textbf{Answer:} if it is "exactly" singular, \mintinline{latex}{SingularMatrixException}
\\~\\
If is near-singular, garbage out \\
So make singularity threshold configurable:
\begin{minted}[xleftmargin=10pt,fontsize=\scriptsize,linenos=false]{java}
OLSMultipleLinearRegression regression =
     new OLSMultipleLinearRegression(threshHold);
\end{minted}
What exactly is this threshHold? \\
What should exception error message say? \\
How to doc it? \\
What if we change underlying impl?
\end{small}
\end{frame}

%-------------- slide -----------------------------------------------%
\begin{frame}[fragile]
  \frametitle{Example 7 - Rank correlation}

Given two double arrays \mintinline{latex}{x[]} and \mintinline{latex}{y[]}, we want
a statistical measure of how well their rank orderings match. \\ ~ \\
Use Spearman's rank correlation coefficient.

\begin{minted}[xleftmargin=10pt,fontsize=\scriptsize,linenos=false]{java}
double[] x = ...
double[] y = ...
SpearmansCorrelation corr = new SpearmansCorrelation();
corr.correlation(x, y);
\end{minted}

How are ties in the data handled? \\
What if one of the arrays contains NaNs?
\end{frame}

%-------------- slide -----------------------------------------------%
\begin{frame}
  \frametitle{Example 7 - Rank correlation (cont)}

\begin{small}
\begin{itemize}
\item As of 3.4.1, by default ties are averaged and the presence of NaN causes RTE
\item Both are configurable, by supplying a RankingAlgorithm to the SpearmansCorrelation
constructor
\item RankingAlgorithm consists of a ranking implementation with configurable
\begin{itemize}
  \item TiesStrategy - what ranks to assign to tied values
  \item NaNStrategy - what to do with NaN values
\end{itemize}
\item Default in Spearman's is natural ranking on doubles with ties averaged and NaNs disallowed
\end{itemize}
Problems: 
\begin{enumerate}
\item What if user gives Spearman's constructor a NaturalRanking with the NaN strategy
that says remove NaNs? 
\item How to doc clearly and completely?
\end{enumerate}
\end{small}
\end{frame}

%----------- slide --------------------------------------------------%
\begin{frame}[fragile]
  \frametitle{Example 8 - Random Numbers}
  
\begin{small}
\begin{itemize}  
\item Commons Math provides alternatives to the JDK PRNG and a pluggable framework
\begin{itemize}
\item RandomGenerator interface is like j.u.Random
\item All random data generation within Commons Math is pluggable
\end {itemize}
\item Well generator is usually the default
\end{itemize}

Low-level example:

\begin{minted}[xleftmargin=10pt,fontsize=\scriptsize,linenos=false]{java}
// Allocate an array to hold random bytes
byte[] bytes = new byte[20];

// Instantiate a Well generator with seed = 100;
Well19937c gen = new Well19937c(100);

// Fill the array with random bytes from the generator
gen.nextBytes(bytes);
\end{minted}
\end{small}
\end{frame}

%----------- slide --------------------------------------------------%
\begin{frame}[fragile]
  \frametitle{Example 8 - Random Numbers (cont)}  
  
High-level example:

\begin{minted}[xleftmargin=10pt,fontsize=\scriptsize,linenos=false]{java}
 // Create a RandomDataGenerator instance using a Mersenne Twister 
 // Seeded with the current time plus system identity hashcode of twister
 RandomDataGenerator gen = 
      new RandomDataGenerator(new MersenneTwister());
      
 // Sample a random value from the Poisson(2) distribution
 long dev = gen.nextPoisson(2);
 \end{minted}

Problems:
\begin{enumerate}
\item Well generators (used as defaults a lot) have some initialization overhead.  How
/ when to take this hit
\item Distribution classes also have sample() methods, but that makes them dependent on
generators
\end{enumerate}
\end{frame}

%----TODO: Nice summary slide reinforcing themes from problems-------%
%--------------------------------------------------------------------%
\section[Getting involved]{Get Involved!}
%----------- slide --------------------------------------------------%
\begin{frame}[fragile]
  \frametitle{Using Apache Commons Math}

Maven:

\begin{minted}[xleftmargin=10pt,fontsize=\small,linenos=false]{xml}
<dependency>
	<groupId>org.apache.commons</groupId>
	<artifactId>commons-math3</artifactId>
	<version>3.4.1</version>
</dependency>
\end{minted}

Download:\\
\begin{small}
\url{http://commons.apache.org/math/download_math.cgi} \\
\end{small}
\\
Watch for new versions!           
\end{frame}

%----------- slide --------------------------------------------------%
\begin{frame}
  \frametitle{Links}

\begin{itemize}
  \item Project homepage: \url{http://commons.apache.org/math/}
  \item Issue tracker: \url{https://issues.apache.org/jira/browse/MATH}
  \item Mailinglists: dev@commons.apache.org \& user@commons.apache.org \\
  e-mail subject: [math]
  \item Wiki: \url{http://wiki.apache.org/commons/MathWishList}
\end{itemize}

\end{frame}

%----------- slide --------------------------------------------------%
\begin{frame}
  \frametitle{How to contribute?}

\begin{itemize}
  \item Check out the user guide \url{http://commons.apache.org/math/userguide/index.html}
  \item Ask questions on the user mailinglist
  \item Participate in discussions on the dev mailinglist
  \item Create bug reports / feature requests / improvements in JIRA
  \item Send patches (code, documentation, examples)
  \item Provide feedback - most welcome!
\end{itemize}

\end{frame}

%----------- slide --------------------------------------------------%
\begin{frame}
  \frametitle{}

\begin{itemize}
  \item Questions?
\end{itemize}

\end{frame}

\end{document}
